        These attributes are not case-sensitive but will be converted to
        lower case.
         
        The attribute values will be applied to all pins.
        
        \btabularx
        \hline
        \multicolumn{4}{|l|}{\bf clock mode} \\
        \hline
        frequency    & 3PL     & float         & frequency, Hz \\
        dutycycle    & 3PL     & float         & duty cycle    \\
        differential & netlist & log           & differential  \\
        ibufg        & netlist & log           & buffer type   \\
        loc          & netlist & str or [2]str & locations     \\
        \hline
        \btabularxend
        
        Attributes {\bf frequency} and {\bf dutycycle} specify these parameters
        for clock nets and any associated clock domains. {\bf Note that at
        present dutycycle is not used. It has been included for future use.}

        The following attribute is read-only and is created when inbuilt
        procedure {\bf clock()} is called and can be retrieved from the clock
        mode variable if required by the function call attribute("ids").
        
        \btabularx
        \hline
        \multicolumn{3}{|l|}{\bf clock mode} \\
        \hline
        ids          & 3PL & []str & element IDs       \\
        \hline
        \btabularxend
        
        The following attributes only apply to a clock variable if it is a clock
        input from a pad. If the {\bf loc} attribute is supplied to a clock
        variable an input pad and buffer will be created, otherwise the clock
        variable source is internal to the FPGA.

        {\bf loc} specifies one or two pad locations (FPGA pins) for the clock
        input - two locations if it is differential (+ve, -ve order). For a
        differential clock input this attribute will be an array of two strings.

        {\bf differential} indicates that the clock input is differential - two
        locations must be supplied and the buffer will be an IBUFDS or IBUFGDS
        rather than the default IBUF or IBUFG. This attribute is used internally
        by 3PL for code generation.

        {\bf ibufg} if true specifies that the input buffer must be an IBUFG or
        IBUFGDS instead of the default IBUF or an IBUFDS. This attribute is used
        internally by 3PL for code generation. It is not clear from Xilinx
        documentation that an IBUFG needs to be explicitly specified, or if
        instead this can be inferred from the package pin(s). other software. 

        The following attributes are required for constraint generation and are
        handled at completion of immediate execution by library procedure {\bf
        generateconstraints()} \autorefph{sec:lpgenerateconstraints}. The
        attribute values are single values which will be applied to the pin or
        pins.
        
        \btabularx
        \hline
        \multicolumn{4}{|l|}{\bf clock mode constraint attributes} \\
        \hline
        tnm          & constraint & str & timing group id                       \\
        diff\_term   & constraint & log & differential  $100\Omega$ termination \\
        iostandard   & constraint & str & voltage standard                      \\
        \hline
        \btabularxend

        {\bf tnm} is the name of a timing group in which all state devices
        clocked by this clock are to be included for timing constraint purposes.
        All clocks which drive logic clock inputs (as distinct from inputs to
        buffers of clock managers) will have a period constraint generated by
        generate\_constraints() in xilinx/common.3pl. By default a timing group
        name will be constructed from the clock identifier translated to upper
        case, but if a "tnm" attribute is provided then this string will be used
        instead.        

        {\bf diff\_term} indicates that a $100\Omega$ termination resistance should
        be applied to a differential input.

        {\bf iostandard} is a string specifying one of the allowed I/O standards
        (see Xilinx documentation).

        The following attribute is read-only and is added when target statement
        variables are assigned or evaluated and the clock domain is established
        or verified. It can be retrieved from the clock mode variable if
        required by the function call attribute("valid").
        
        \btabularx
        \hline
        \multicolumn{3}{|l|}{\bf clock mode} \\
        \hline
        valid   & 3PL & log & is a valid logic clock    \\
        \hline
        \btabularxend

        Any attributes that are not listed above will result in a fatal error.
